\section{Distributed Network Topology and Hardware Interface}

To optimize computational load and ensure modular system performance, we implemented a distributed computing architecture leveraging the Data Distribution Service (DDS) inherent to ROS 2. As illustrated in Figure~\ref{fig:dist_arch}, rather than relying on a single processing unit, the system workload is partitioned across two separate laptops, designated as \textbf{Laptop 1} and \textbf{Laptop 2}.

\begin{figure}[H]
    \centering
    \includegraphics[width=0.9\textwidth]{Figures/chapter3/Divided_network/Divided_Computation.png} 
    \caption{Hardware and Network Communication Architecture for the Collaborative System.}
    \label{fig:dist_arch}
\end{figure}
\subsection{Physical Network Layer}
The core communication backbone is established via a dedicated local area network (LAN). Both laptops and the ABB robot controller are connected to a central D-Link Fast Ethernet switch. This star topology ensures low-latency, high-bandwidth communication between the processing units and the industrial manipulators.

\begin{itemize}
    \item \textbf{ABB Robot Integration:} The ABB arm is networked directly to the Ethernet switch. It communicates via \textbf{UDP}, ensuring fast data transmission for real-time feedback and control loops required by the ROS 2 bridge.
    
    \item \textbf{AR4 Robot Integration:} Unlike the networked ABB, the AR4 robot is interfaced directly to \textbf{Laptop 2} via a USB-C connection, utilizing serial communication for direct motor control instructions.
\end{itemize}


\subsection{Distributed ROS 2 Environment}
By configuring both laptops to operate on the same \texttt{ROS\_DOMAIN\_ID}, the ROS2 nodes running on separate machines function as a unified computation graph. This allows topics published by the driver on the AR4-connected laptop to be immediately subscribed to by the planning or perception nodes on the second laptop without complex routing protocols.

This separation allows for significant resource isolation based on task requirements:

\begin{itemize}
    \item \textbf{Laptop 1 (High-Level Computation):} This machine acts as the system supervisor. It handles computationally intensive tasks including the GUI, Object Detection, and Grasping Strategy. It is physically connected to the \textbf{Intel RealSense} camera via USB/Serial to process high-fidelity depth data.
    
    \item \textbf{Laptop 2 (Real-Time Control):} This machine is dedicated to hardware interfacing. It runs the AR4 Drivers, Controllers, and IBVS (Image-Based Visual Servoing) algorithms. Additionally, it handles telemetry from the \textbf{ESP32 Camera}, which connects wirelessly via WiFi to this specific node.
\end{itemize}